\documentclass[12pt,a4paper]{article}

% ---------------- ОПРЕДЕЛЕНИЕ ДВИЖКА ----------------
\usepackage{iftex}

% ---------------- ШРИФТЫ И КОДИРОВКИ ----------------
\ifPDFTeX
  \usepackage[utf8]{inputenc}
  \usepackage[T2A,T1]{fontenc}
  \usepackage{lmodern}
\else
  \usepackage{fontspec}
  \defaultfontfeatures{Ligatures=TeX}   % --- , ``'' работают как в обычном TeX
  % CM Unicode: те же формы, что у Computer Modern, но с кириллицей.
  % Шрифты задаются именами файлов, поэтому не требуют установки в систему.
  \setmainfont{cmunrm.otf}[
    BoldFont       = cmunbx.otf,
    ItalicFont     = cmunti.otf,
    BoldItalicFont = cmunbi.otf]
  \setsansfont{cmunss.otf}[
    BoldFont       = cmunsx.otf,
    ItalicFont     = cmunsi.otf,
    BoldItalicFont = cmunso.otf]
  \setmonofont{cmuntt.otf}[
    BoldFont       = cmuntb.otf,
    ItalicFont     = cmunit.otf,
    BoldItalicFont = cmuntx.otf,
    Mapping        = {}]   % иначе апострофы в verbatim станут типографскими
\fi

\usepackage[russian, english]{babel}   % грузить ПОСЛЕ fontspec/fontenc

% ---------------- МАТЕМАТИКА ----------------
\usepackage{amsmath, amssymb, amsthm, mathtools}

\numberwithin{equation}{section}

\theoremstyle{plain}
\newtheorem{theorem}{Теорема}[section]
\newtheorem{lemma}[theorem]{Лемма}
\newtheorem{proposition}[theorem]{Утверждение}
\theoremstyle{definition}
\newtheorem{definition}[theorem]{Определение}
\theoremstyle{remark}
\newtheorem{remark}[theorem]{Замечание}

% ---------------- ВЁРСТКА ----------------
\usepackage[top=1in, bottom=1in, left=0.5in, right=0.5in]{geometry}
\ifPDFTeX
  \usepackage[protrusion=true, expansion=true]{microtype}
\else
  \usepackage[protrusion=true, expansion=false]{microtype}
\fi
\usepackage{indentfirst}
\usepackage{setspace}

% ---------------- СПИСКИ, ТАБЛИЦЫ, РИСУНКИ ----------------
\usepackage{enumitem}
\usepackage{booktabs}
\usepackage{float}
\usepackage{graphicx}
\usepackage{subcaption}

% ---------------- ОГЛАВЛЕНИЕ ----------------
\usepackage{tocloft}
\renewcommand{\cftsecpagefont}{\bfseries}

% ---------------- ССЫЛКИ ----------------
\usepackage[unicode=true, colorlinks=true,
            linkcolor=blue, urlcolor=blue, citecolor=blue]{hyperref}

\title{Эконометрика. Лекции и семинары}
\author{Баулин Евгений}
\date{\today}

\begin{document}
\selectlanguage{russian}

\maketitle
\tableofcontents
\clearpage

\section{Вводная часть. 7 сентября 2026}

Лектора зовут Анна Жимерикина

Telegram: \href{https://t.me/ann_zhi}{\texttt{@ann\_zhi}}

Email: \href{mailto:aj.zhimerikina@physics.msu.ru}{\nolinkurl{aj.zhimerikina@physics.msu.ru}}

О всех изменениях о парах будут сообщать в ТГК: \href{https://t.me/ftiad_econ}{\texttt{@ftiad\_econ}}

Правила посещения:

Отметок о посещении нет, поэтому посещение свободное (можно приходить, а можно и не приходить)

Тесты на семинарах как доп к оценке, то есть не писав их всё ещё можно получить $10$. Тесты дадут максимум $1-1.5$

Ровно в 9 pm заканчивается

Программировать будем только на Python

$$
    \text{Итоговая оценка} = 0.3 * \text{ДЗ}_1 + 0.4 * \text{ДЗ}_2 + 0.3 * \text{КР}
$$

Можно просто написать в ЛС лектору, чтобы получить 4 за курс и не ходить на пары

Материалы: \href{https://drive.google.com/drive/folders/1e8s9vTZWyk2HHJh6tHfEJ-FtTs13sTG0}{\texttt{Google Drive}}

Ноутбук — опционально, можно не приносить, будут показывать экран

\section{Лекция и семинар. 7 сентября 2026}

\subsection{Зачем нужна эконометрика, когда есть продвинутое машинное обучение}

\begin{itemize}
    \item Когда данных очень мало, а результат нужно получить сейчас
    \item Когда нужно полностью интерпретировать модель
    \item Всякие экономические модели
\end{itemize}

\subsection{Типы данных в эконометрике}

\begin{itemize}
    \item Временные ряды (time-series data)
    \item Пространственные данные (cross-sectional data)
    \item Панельные данные (panel data)
\end{itemize}

Примеры смотреть в файле Лекция 1.pdf

\subsection{Глоссарий}

Всё, что $y$ — target, целевая переменная, зависимая переменная, объясняемая переменная

Остальное — регрессоры, объясняющие переменные, features, признаки

По каждой переменной $n$ наблюдений $y_1, y_2, \dots, y_n$

\subsection{Пример эконометрической задачи}

Как скорость машины влияет на длину тормозного пути? На каком расстоянии от пешеходного перехода установить знак?

Данные: Ezekiel, M. (1930) Methods of Correlation Analysis. Wiley (классический датасет \texttt{cars}: скорость в милях в час и длина тормозного пути в футах)

График смотреть в файле Лекция 1.pdf

\subsection{Модель}

Пусть мы знаем, что в реальности данные связаны следующим образом:

$$
    y_i = \beta_1 + \beta_2 x_i + \varepsilon_i
$$

\begin{itemize}
    \item Наблюдаемые нами переменные (наш датасет): $y, x$
    \item Неизвестные параметры: $\beta_1, \beta_2$. Они существуют, но наблюдать мы их не можем
    \item Случайная составляющая, ошибка: $\varepsilon$
\end{itemize}

Ошибка возникает из-за несовершенства реальности:

\begin{itemize}
    \item Ошибки измерений
    \item Неопределённость человеческого поведения
    \item Неучёт влияния бесчисленного множества случайных событий
\end{itemize}

Чтобы строить модели нужно включить здравый смысл и понять, что ошибка в измерениях~— это нормально. Не нужно добавлять всё, что угодно (цену курицы, цвет машины и т.\,п.)

$\beta_1 + \beta_2 x_i$ — детерминированная часть

$\varepsilon_i$ — стохастическая часть

\subsection{Задача эконометрики}

План действий:

\begin{enumerate}
    \item Придумать адекватную модель (так как реальное устройство мира нам неизвестно, мы можем лишь делать предположение)
    \item Получить оценки неизвестных параметров: $\hat{\beta_1}, \hat{\beta_2}$
    \item Прогнозировать, заменив неизвестные параметры на оценки: $\hat{y_i} = \hat{\beta_1} + \hat{\beta_2}x_i$
\end{enumerate}

\subsection{Метод наименьших квадратов (МНК)}

Нужен способ получить оценки неизвестных параметров модели исходя из реальных данных

Ошибка прогноза: $\hat{e_i} = y_i - \hat{y_i}$

Сумма квадратов ошибок прогноза:
$$
    Q \left(\hat{\beta_1}, \hat{\beta_2}\right) = \sum_{i = 1}^n \hat{e_i}^2 = \sum_{i = 1}^n \left(y_i - \hat{y_i}\right)^2
$$

Суть МНК: В качестве оценок взять такие $\hat{\beta_1}, \hat{\beta_2}$, при которых сумма квадратов ошибок прогноза $Q$ минимальна

\subsection{Частные случаи МНК}

\subsubsection{Регрессия на константу}

В модели $y_i = \beta + \varepsilon_i$: $\hat{\beta} = \bar{y}$

Интерпретация: В модели без объясняющих переменных наилучший прогноз — это среднее значение зависимой переменной

\subsubsection{Парная регрессия}

Модель с константой и одной независимой переменной

В модели $y_i = \beta_1 + \beta_2 x_i + \varepsilon_i$

\begin{gather*}
    \hat{\beta_2} = \frac{\sum \left(x_i - \bar{x}\right)\left(y_i - \bar{y}\right)}{\sum \left(x_i - \bar{x}\right)^2} = \frac{s Cov (x, y)}{s Var(x)} \\
    \hat{\beta_1} = \bar{y} - \hat{\beta_2}\bar{x}
\end{gather*}

Замечательное свойство: Точка $\left(\bar{x}, \bar{y}\right)$ лежит на линии регрессии $\hat{y_i} = \hat{\beta_1} + \hat{\beta_2} x$

\subsection{Терминология и обозначения}

\begin{itemize}
    \item $y_i$ — таргет, зависимая, объясняемая переменная
    \item $x_i$ — feature, регрессор, независимая, объясняющая переменная
    \item $\varepsilon$ - ошибка (несовершенство мира), случайная составляющая
    \item $\hat{y_i}$ — прогноз, прогнозное значение
    \item $\hat{e_i} = y_i - \hat{y_i}$ — остаток, ошибка прогноза (не путайте с $\varepsilon$)
    \item $RSS = \sum_{i = 1}^{n} \hat{e_i}^2$ — сумма квадратов остатков
\end{itemize}

\subsection{Случай множественной регрессии}

Регрессоров может быть несколько, поэтому модель удобно записывать в матричном виде

Вектор зависимой переменной, вектор прогнозных значений и вектор случайных ошибок:
$$
    y = \begin{pmatrix} y_1 \\ y_2 \\ \vdots \\ y_n \end{pmatrix}, \qquad
    \hat{y} = \begin{pmatrix} \hat{y_1} \\ \hat{y_2} \\ \vdots \\ \hat{y_n} \end{pmatrix}, \qquad
    \varepsilon = \begin{pmatrix} \varepsilon_1 \\ \varepsilon_2 \\ \vdots \\ \varepsilon_n \end{pmatrix}
$$

Матрица регрессоров (первый столбец — константа), вектор коэффициентов и вектор их оценок:
$$
    X = \begin{pmatrix}
        1      & x_1    & z_1    \\
        1      & x_2    & z_2    \\
        \vdots & \vdots & \vdots \\
        1      & x_n    & z_n
    \end{pmatrix}, \qquad
    \beta = \begin{pmatrix} \beta_1 \\ \beta_2 \\ \beta_3 \end{pmatrix}, \qquad
    \hat{\beta} = \begin{pmatrix} \hat{\beta_1} \\ \hat{\beta_2} \\ \hat{\beta_3} \end{pmatrix}
$$

Спецификация модели:
$$
    y = X\beta + \varepsilon
$$

\subsubsection{Минимизация суммы квадратов остатков}

Вектор остатков: $e = y - \hat{y} = y - X\hat{\beta}$

Сумма квадратов остатков: $RSS = \sum \hat{e_i}^2 = e'e \to \min$

Выразим $e'e$ через $X$ и $\hat{\beta}$:
\begin{gather*}
    e'e = \left(y - X\hat{\beta}\right)'\left(y - X\hat{\beta}\right) = \\
    = y'y - y'X\hat{\beta} - \hat{\beta}'X'y + \hat{\beta}'X'X\hat{\beta} = \\
    = y'y - 2\hat{\beta}'X'y + \hat{\beta}'X'X\hat{\beta}
\end{gather*}

Слагаемые $y'X\hat{\beta}$ и $\hat{\beta}'X'y$ — это числа (матрицы размера $1 \times 1$), причём одно транспонировано к другому, поэтому они равны и складываются в $2\hat{\beta}'X'y$

\subsubsection{Оценка метода наименьших квадратов}

Необходимые условия минимума $RSS$ получаются дифференцированием по вектору $\hat{\beta}$:
$$
    \frac{\partial RSS}{\partial \hat{\beta}} = -2X'y + 2X'X\hat{\beta} = 0
$$

Решая уравнение, получаем оценку метода наименьших квадратов:
$$
    \hat{\beta}_{OLS} = \left(X'X\right)^{-1}X'y
$$

Предположение: матрица $X'X$ обратима (столбцы $X$ линейно независимы)

Если это не так (полная мультиколлинеарность), то оценка МНК не единственна

\subsection{Суммы квадратов}

\begin{itemize}
    \item Сумма квадратов остатков, $RSS$ (residual sum of squares):
          $$
              RSS = \sum \hat{e_i}^2
          $$
    \item Общая сумма квадратов, $TSS$ (total sum of squares) — вся дисперсия:
          $$
              TSS = \sum \left(y_i - \bar{y}\right)^2
          $$
    \item Объяснённая сумма квадратов, $ESS$ (explained sum of squares):
          $$
              ESS = \sum \left(\hat{y_i} - \bar{y}\right)^2
          $$
\end{itemize}

Интерпретация $ESS$: $ESS$ измеряет, насколько хорошо модель объяснила вариацию $y$ по сравнению с простым средним $\bar{y}$. Чем больше доля $ESS$ в $TSS$, тем лучше модель описывает данные

\subsection{Если в регрессию включён свободный член}

Если в регрессию включён свободный член, $y_i = \beta_1 + \dots$, и оценки МНК единственны, то:

\begin{itemize}
    \item $\sum \hat{e_i} = 0$
    \item $\sum y_i = \sum \hat{y_i}$
    \item $\bar{y} = \overline{\hat{y}}$
    \item $TSS = RSS + ESS$
\end{itemize}

$$
    \underbrace{\sum \left(y_i - \bar{y}\right)^2}_{TSS} = \underbrace{\sum \left(y_i - \hat{y_i}\right)^2}_{RSS} + \underbrace{\sum \left(\hat{y_i} - \bar{y}\right)^2}_{ESS}
$$

\subsection{Коэффициент детерминации — простой показатель качества}

Коэффициентом детерминации, или долей объяснённой дисперсии, называется
$$
    R^2 = 1 - \frac{RSS}{TSS} = \frac{ESS}{TSS}
$$

В силу определения $R^2$ принимает значения между $0$ и $1$: $0 \leqslant R^2 \leqslant 1$

\begin{itemize}
    \item $R^2 = 0$ означает, что регрессия ничего не даёт, то есть $x_i$ не улучшает качество предсказания $y_i$ по сравнению с тривиальным предсказанием $\hat{y_i} = \bar{y}$
    \item $R^2 = 1$ означает точную подгонку: все точки наблюдений лежат на регрессионной прямой (все $\hat{e_i} = 0$)
\end{itemize}

Важное замечание: $0 \leqslant R^2 \leqslant 1$ выполняется только при прогнозировании in sample. При прогнозе out-of-sample $R^2$ может принимать значения из $(-\infty; 1]$, то есть прогнозировать хуже выборочного среднего $\bar{y}$

Теорема: если в регрессию включён свободный член, $y_i = \beta_1 + \dots$, и оценки МНК единственны, то $R^2$ равен квадрату выборочной корреляции между $y$ и $\hat{y}$:
$$
    R^2 = \left(sCorr\left(y, \hat{y}\right)\right)^2 = \left(\frac{\sum \left(y_i - \bar{y}\right)\left(\hat{y_i} - \bar{y}\right)}{\sqrt{\sum \left(y_i - \bar{y}\right)^2} \sqrt{\sum \left(\hat{y_i} - \bar{y}\right)^2}}\right)^2
$$

\subsubsection{Коэффициент детерминации в общем случае}

$$
    R^2 \left(y, \hat{y}\right) = 1 - \frac{\sum_{i = 1}^n \left(y_i - \hat{y_i}\right)^2}{\sum_{i = 1}^n \left(y_i - \bar{y}\right)^2} = 1 - \frac{MSE}{Var(y)}
$$

\begin{itemize}
    \item $MSE$ может быть бесконечно большой величиной
    \item В библиотеке \texttt{sklearn} эта формула реализована в \texttt{sklearn.metrics.r2\_score}
\end{itemize}

Из документации \texttt{sklearn}: функция не симметрична, и, в отличие от большинства других метрик, $R^2$ может быть отрицательным (он не обязан быть квадратом какой-либо величины $R$)

\subsection{Разбор с семинара}

На семинаре выводим формулы, которые на лекции были даны в готовом виде, и считаем всё то же самое руками и на Python

\subsubsection{Что нужно помнить из теории вероятностей и статистики}

Все формулы МНК записываются через выборочные моменты. Для выборок $x_1, \dots, x_n$ и $y_1, \dots, y_n$:

\begin{itemize}
    \item Выборочное среднее: $\bar{x} = \frac{1}{n}\sum_{i = 1}^n x_i$
    \item Выборочная дисперсия: $sVar(x) = \frac{1}{n}\sum_{i = 1}^n \left(x_i - \bar{x}\right)^2$
    \item Выборочная ковариация: $sCov(x, y) = \frac{1}{n}\sum_{i = 1}^n \left(x_i - \bar{x}\right)\left(y_i - \bar{y}\right)$
    \item Выборочная корреляция: $sCorr(x, y) = \dfrac{sCov(x, y)}{\sqrt{sVar(x)}\sqrt{sVar(y)}}$
\end{itemize}

Замечание: в формуле для $\hat{\beta_2}$ множители $\frac{1}{n}$ в числителе и знаменателе сокращаются, поэтому не важно, делить суммы на $n$ или на $n - 1$ — оценка коэффициента от этого не изменится

Тождества, которые постоянно используются в выкладках:

\begin{itemize}
    \item $\sum \left(x_i - \bar{x}\right) = 0$
    \item $\sum \left(x_i - \bar{x}\right)\left(y_i - \bar{y}\right) = \sum \left(x_i - \bar{x}\right) y_i = \sum x_i y_i - n \bar{x}\bar{y}$
    \item $\sum \left(x_i - \bar{x}\right)^2 = \sum x_i \left(x_i - \bar{x}\right) = \sum x_i^2 - n\bar{x}^2$
\end{itemize}

\subsubsection{Вывод МНК-оценки: регрессия на константу}

Модель $y_i = \beta + \varepsilon_i$, прогноз $\hat{y_i} = \hat{\beta}$. Минимизируем
$$
    Q\left(\hat{\beta}\right) = \sum_{i = 1}^n \left(y_i - \hat{\beta}\right)^2
$$

Необходимое условие минимума:
$$
    \frac{dQ}{d\hat{\beta}} = -2\sum_{i = 1}^n \left(y_i - \hat{\beta}\right) = 0
    \quad \Longrightarrow \quad
    \sum_{i = 1}^n y_i = n\hat{\beta}
    \quad \Longrightarrow \quad
    \hat{\beta} = \bar{y}
$$

Это действительно минимум: $\dfrac{d^2 Q}{d\hat{\beta}^2} = 2n > 0$, то есть $Q$ строго выпукла

\subsubsection{Вывод МНК-оценок: парная регрессия}

Модель $y_i = \beta_1 + \beta_2 x_i + \varepsilon_i$, прогноз $\hat{y_i} = \hat{\beta_1} + \hat{\beta_2} x_i$. Минимизируем
$$
    Q\left(\hat{\beta_1}, \hat{\beta_2}\right) = \sum_{i = 1}^n \left(y_i - \hat{\beta_1} - \hat{\beta_2} x_i\right)^2
$$

Приравниваем к нулю частные производные и получаем систему нормальных уравнений:
\begin{gather*}
    \frac{\partial Q}{\partial \hat{\beta_1}} = -2\sum \left(y_i - \hat{\beta_1} - \hat{\beta_2} x_i\right) = 0 \\
    \frac{\partial Q}{\partial \hat{\beta_2}} = -2\sum x_i \left(y_i - \hat{\beta_1} - \hat{\beta_2} x_i\right) = 0
\end{gather*}

Из первого уравнения:
$$
    \sum y_i = n\hat{\beta_1} + \hat{\beta_2} \sum x_i
    \quad \Longrightarrow \quad
    \bar{y} = \hat{\beta_1} + \hat{\beta_2}\bar{x}
    \quad \Longrightarrow \quad
    \hat{\beta_1} = \bar{y} - \hat{\beta_2}\bar{x}
$$

Это ровно то самое замечательное свойство с лекции: точка $\left(\bar{x}, \bar{y}\right)$ лежит на линии регрессии

Подставляем $\hat{\beta_1}$ во второе уравнение:
\begin{gather*}
    \sum x_i \left(y_i - \bar{y} + \hat{\beta_2}\bar{x} - \hat{\beta_2} x_i\right) = 0 \\
    \sum x_i \left(y_i - \bar{y}\right) = \hat{\beta_2} \sum x_i \left(x_i - \bar{x}\right)
\end{gather*}

Пользуясь тождествами выше, получаем формулу с лекции:
$$
    \hat{\beta_2} = \frac{\sum \left(x_i - \bar{x}\right)\left(y_i - \bar{y}\right)}{\sum \left(x_i - \bar{x}\right)^2} = \frac{sCov(x, y)}{sVar(x)}
$$

Интерпретация: $\hat{\beta_2}$ показывает, на сколько единиц в среднем изменится $y$, если $x$ вырастет на одну единицу. Знак $\hat{\beta_2}$ совпадает со знаком выборочной ковариации $x$ и $y$

\subsubsection{Свойства остатков МНК}

Нормальные уравнения — это в точности условия ортогональности остатков регрессорам. Отсюда сразу следуют свойства, которые на лекции были даны списком:

\begin{itemize}
    \item $\sum \hat{e_i} = 0$ — первое нормальное уравнение (верно, только если в модели есть константа)
    \item $\sum x_i \hat{e_i} = 0$ — второе нормальное уравнение
    \item $\bar{y} = \overline{\hat{y}}$ — следствие первого: $\sum \hat{y_i} = \sum \left(y_i - \hat{e_i}\right) = \sum y_i$
\end{itemize}

В матричном виде все нормальные уравнения записываются одной строкой:
$$
    X'e = 0
    \quad \Longleftrightarrow \quad
    X'\left(y - X\hat{\beta}\right) = 0
    \quad \Longleftrightarrow \quad
    X'X\hat{\beta} = X'y
$$

Если матрица $X'X$ обратима, отсюда получается формула с лекции: $\hat{\beta}_{OLS} = \left(X'X\right)^{-1}X'y$

Геометрический смысл: $\hat{y} = X\left(X'X\right)^{-1}X'y$ — проекция вектора $y$ на линейную оболочку столбцов $X$, а вектор остатков $e$ ортогонален этой оболочке. МНК ищет ближайшую к $y$ точку среди тех, которые модель вообще способна описать

\subsubsection{Разложение общей суммы квадратов}

Запишем отклонение от среднего как сумму двух слагаемых:
$$
    y_i - \bar{y} = \left(y_i - \hat{y_i}\right) + \left(\hat{y_i} - \bar{y}\right) = \hat{e_i} + \left(\hat{y_i} - \bar{y}\right)
$$

Возводим в квадрат и суммируем:
$$
    \sum \left(y_i - \bar{y}\right)^2 = \sum \hat{e_i}^2 + \sum \left(\hat{y_i} - \bar{y}\right)^2 + 2\sum \hat{e_i}\left(\hat{y_i} - \bar{y}\right)
$$

Осталось показать, что перекрёстное слагаемое равно нулю:
$$
    \sum \hat{e_i}\left(\hat{y_i} - \bar{y}\right) = \sum \hat{e_i}\left(\hat{\beta_1} + \hat{\beta_2}x_i - \bar{y}\right) = \left(\hat{\beta_1} - \bar{y}\right)\underbrace{\sum \hat{e_i}}_{= \, 0} + \hat{\beta_2}\underbrace{\sum x_i \hat{e_i}}_{= \, 0} = 0
$$

Получаем $TSS = RSS + ESS$. Обратите внимание: доказательство целиком опирается на нормальные уравнения, поэтому без константы в модели разложение, вообще говоря, ломается

\subsubsection{Коэффициент детерминации в парной регрессии}

Задача: показать, что в парной регрессии $R^2 = \left(sCorr(x, y)\right)^2$

Решение. Так как $\hat{y_i} - \bar{y} = \hat{\beta_2}\left(x_i - \bar{x}\right)$, то
$$
    ESS = \hat{\beta_2}^2 \sum \left(x_i - \bar{x}\right)^2 = \frac{\left(\sum \left(x_i - \bar{x}\right)\left(y_i - \bar{y}\right)\right)^2}{\sum \left(x_i - \bar{x}\right)^2}
$$

Делим на $TSS = \sum \left(y_i - \bar{y}\right)^2$:
$$
    R^2 = \frac{ESS}{TSS} = \frac{\left(\sum \left(x_i - \bar{x}\right)\left(y_i - \bar{y}\right)\right)^2}{\sum \left(x_i - \bar{x}\right)^2 \sum \left(y_i - \bar{y}\right)^2} = \left(sCorr(x, y)\right)^2
$$

То есть в парной регрессии $R^2$ не различает, что на что регрессируем: он одинаков для регрессии $y$ на $x$ и для регрессии $x$ на $y$

\subsubsection{Числовой пример}

\begin{table}[H]
    \centering
    \begin{tabular}{lrrrrr}
        \toprule
        $x_i$ & 1 & 2 & 3 & 4 & 5 \\
        $y_i$ & 2 & 4 & 5 & 4 & 5 \\
        \bottomrule
    \end{tabular}
\end{table}

Считаем средние: $\bar{x} = 3$, $\bar{y} = 4$

\begin{gather*}
    \sum \left(x_i - \bar{x}\right)\left(y_i - \bar{y}\right) = (-2)(-2) + (-1) \cdot 0 + 0 \cdot 1 + 1 \cdot 0 + 2 \cdot 1 = 6 \\
    \sum \left(x_i - \bar{x}\right)^2 = 4 + 1 + 0 + 1 + 4 = 10
\end{gather*}

Отсюда
$$
    \hat{\beta_2} = \frac{6}{10} = 0.6, \qquad
    \hat{\beta_1} = 4 - 0.6 \cdot 3 = 2.2, \qquad
    \hat{y_i} = 2.2 + 0.6 x_i
$$

Прогнозы и остатки:

\begin{table}[H]
    \centering
    \begin{tabular}{lrrrrr}
        \toprule
        $\hat{y_i}$ & 2.8    & 3.4 & 4.0 & 4.6    & 5.2    \\
        $\hat{e_i}$ & $-0.8$ & 0.6 & 1.0 & $-0.6$ & $-0.2$ \\
        \bottomrule
    \end{tabular}
\end{table}

Проверка свойств остатков: $\sum \hat{e_i} = 0$ и $\sum x_i \hat{e_i} = 0$

Суммы квадратов: $RSS = 2.4$, $ESS = \hat{\beta_2}^2 \cdot 10 = 3.6$, $TSS = 6$. Действительно, $TSS = RSS + ESS$, а
$$
    R^2 = 1 - \frac{2.4}{6} = 0.6
$$

\subsubsection{Как это считается на Python}

\begin{verbatim}
import numpy as np
import pandas as pd
import statsmodels.api as sm

df = pd.read_csv('cars.csv')   # speed (мили в час), dist (футы)
x = df['speed'].to_numpy()
y = df['dist'].to_numpy()

# 1. МНК руками по формулам парной регрессии
beta2 = np.sum((x - x.mean()) * (y - y.mean())) / np.sum((x - x.mean()) ** 2)
beta1 = y.mean() - beta2 * x.mean()

# 2. МНК в матричном виде
X = sm.add_constant(x)                      # добавляем столбец из единиц
beta = np.linalg.solve(X.T @ X, X.T @ y)    # надёжнее, чем np.linalg.inv

# 3. МНК через statsmodels
model = sm.OLS(y, X).fit()
print(model.params)          # те же beta1, beta2
print(model.summary())

# 4. Суммы квадратов и коэффициент детерминации
y_hat = X @ beta
e = y - y_hat
rss = np.sum(e ** 2)
tss = np.sum((y - y.mean()) ** 2)
ess = np.sum((y_hat - y.mean()) ** 2)

print(rss + ess, tss)                 # совпадают
print(1 - rss / tss, model.rsquared)  # тоже совпадают
print(e.sum(), (x * e).sum())         # оба нуля с точностью до машинной
\end{verbatim}

Полезно помнить:

\begin{itemize}
    \item Константа не добавляется сама: без \texttt{sm.add\_constant} модель оценивается без свободного члена, и ломаются и $R^2$, и свойства остатков
    \item Считать $\left(X'X\right)^{-1}$ через \texttt{np.linalg.inv} не нужно: \texttt{np.linalg.solve} и \texttt{np.linalg.lstsq} точнее и устойчивее численно
    \item В \texttt{sklearn} за константу в \texttt{LinearRegression} отвечает параметр \texttt{fit\_intercept} (по умолчанию \texttt{True}), а $R^2$ считает метод \texttt{score}
\end{itemize}

\subsubsection{Подводные камни}

\begin{itemize}
    \item $R^2$ не убывает при добавлении любого нового регрессора, даже бессмысленного, поэтому сравнивать по нему модели с разным числом регрессоров нельзя
    \item Регрессия $y$ на $x$ и регрессия $x$ на $y$ дают разные прямые. Их коэффициенты связаны соотношением $\hat{\beta}_{y|x} \cdot \hat{\beta}_{x|y} = R^2$
    \item Высокий $R^2$ ничего не говорит о причинно-следственной связи
    \item Единицы измерения: если заменить $x$ на $cx$, то $\hat{\beta_2}$ поделится на $c$, а прогнозы, остатки и $R^2$ не изменятся
\end{itemize}

\subsubsection{Задачи для самостоятельного решения}

\begin{enumerate}
    \item Показать, что в модели без константы $y_i = \beta x_i + \varepsilon_i$ МНК-оценка равна $\hat{\beta} = \dfrac{\sum x_i y_i}{\sum x_i^2}$
    \item Показать, что в модели без константы, вообще говоря, $\sum \hat{e_i} \neq 0$, разложение $TSS = RSS + ESS$ не выполняется, а $R^2$, посчитанный по формуле $1 - RSS / TSS$, может оказаться отрицательным
    \item Все $y_i$ увеличили на константу $c$. Как изменятся $\hat{\beta_1}$, $\hat{\beta_2}$, $RSS$ и $R^2$?
    \item Показать, что $\hat{\beta_2} = sCorr(x, y) \cdot \sqrt{\dfrac{sVar(y)}{sVar(x)}}$
    \item Взять датасет \texttt{cars} из лекции, построить регрессию длины тормозного пути на скорость и проинтерпретировать оба коэффициента. Имеет ли содержательный смысл $\hat{\beta_1}$?
\end{enumerate}

\end{document}
